\documentclass[12pt,a4paper]{article}

% ---------- Pacotes ----------
\usepackage[utf8]{inputenc}
\usepackage[T1]{fontenc}
\usepackage[brazilian]{babel}
\usepackage[a4paper,top=3cm,bottom=2cm,left=3cm,right=2cm]{geometry}
\usepackage{setspace}
\usepackage{indentfirst}
\usepackage{fancyhdr}
\usepackage{url}
\usepackage{hyperref}

\hypersetup{
    hidelinks,
    pdftitle={Resenha - Ciência e ideologia},
    pdfauthor={Jeann Victor Batista},
    pdfsubject={Filosofia e Metodologia da Ciência}
}
\setstretch{1.5}
\setlength{\parindent}{1.25cm}
\setlength{\parskip}{0pt}
\setlength{\emergencystretch}{3em}

\pagestyle{fancy}
\fancyhf{}
\fancyfoot[C]{\thepage}
\renewcommand{\headrulewidth}{0pt}

% ---------- Dados do discente ----------
\newcommand{\NomeAluno}{Jeann Victor Batista}
\newcommand{\InstituicaoCurso}{Universidade Federal de Alfenas | Ciência da Computação}
\newcommand{\DisciplinaProfessor}{Filosofia e Metodologia da Ciência - José Francisco Lopes Xarão}

\begin{document}

% ======================================================
% CAPA / IDENTIFICAÇÃO
% ======================================================
\begin{center}
    \textbf{\NomeAluno}\footnote{Em conformidade com a Resolução nº 02/2025 (Diretrizes para Uso de Inteligência Artificial na UNIFAL-MG), declara-se que houve uso de ferramenta de Inteligência Artificial da OpenAI como apoio na organização, adequação ao modelo e revisão linguística do texto. A leitura da obra, a interpretação do conteúdo e a responsabilidade pela versão entregue permanecem com o discente.}\\[2pt]
    \InstituicaoCurso\\[2pt]
    \DisciplinaProfessor\\[24pt]

    \textbf{\Large RESENHA}\\[6pt]
    \textit{Ciência e ideologia}\\[4pt]
    \textit{\small Science and ideology}
\end{center}

\vspace{18pt}

% ======================================================
% REFERÊNCIA DA OBRA
% ======================================================
\noindent\textbf{Resenha da obra:}\par
\noindent\setlength{\parindent}{0pt}%
FOUREZ, Gérard. Ciência e ideologia. In: FOUREZ, Gérard. \textit{A construção das ciências: introdução à filosofia e à ética das ciências}. Tradução de Luiz Paulo Rouanet. São Paulo: Editora da Universidade Estadual Paulista, 1995. cap. 7, p. 179--193.
\setlength{\parindent}{1.25cm}

\vspace{12pt}

% ======================================================
% RESUMO TÉCNICO-CIENTÍFICO
% ======================================================
\noindent\textbf{Resumo}\par
\noindent\setlength{\parindent}{0pt}%
Esta resenha examina o capítulo ``Ciência e ideologia'', de Gérard Fourez, publicado no livro \textit{A construção das ciências: introdução à filosofia e à ética das ciências}. O tema do capítulo é a relação contraditória entre ciência e ideologia. O problema discutido pelo autor é saber como a ciência pode criticar discursos ideológicos se ela própria é uma construção histórica, orientada por paradigmas, interesses e escolhas. Fourez defende que os testes científicos permitem questionar afirmações gerais e revelar alguns de seus limites, mas não oferecem uma visão neutra e completa da realidade. Para sustentar essa posição, diferencia ideologias de primeiro e segundo graus, examina a incapacidade da ciência de resolver sozinha questões éticas e analisa a influência dos grupos sociais sobre a produção do conhecimento. A exposição é conceitual e se apoia em exemplos cotidianos e científicos. Na apreciação crítica, destaca-se como mérito a recusa tanto da neutralidade absoluta da ciência quanto da ideia de que todo conhecimento seria arbitrário. Aponta-se, porém, que a distinção entre os graus de ideologia depende de um julgamento cujo critério poderia ser mais desenvolvido. A discussão é relacionada ao filme \textit{Oppenheimer}, no qual o domínio técnico para produzir a bomba não resolve o problema moral de seu uso.

\vspace{6pt}
\noindent\textbf{Palavras-chave:} Ciência. Ideologia. Ética. Discurso científico. Conhecimento.
\setlength{\parindent}{1.25cm}

\vspace{12pt}

% ======================================================
% ABSTRACT
% ======================================================
\noindent\textbf{Abstract}\par
\noindent\setlength{\parindent}{0pt}%
\textit{This review examines the chapter ``Science and Ideology'', by Gérard Fourez, published in the book A construção das ciências: introdução à filosofia e à ética das ciências. The chapter addresses the contradictory relationship between science and ideology. Its main problem is how science can criticize ideological discourses if scientific knowledge is itself a historical construction guided by paradigms, interests, and choices. Fourez argues that scientific tests can challenge broad claims and reveal some of their limits, but cannot provide a neutral and complete view of reality. He distinguishes first- and second-degree ideologies, discusses the inability of science to solve ethical questions by itself, and examines the influence of social groups on knowledge production. His conceptual discussion relies on everyday and scientific examples. The critical assessment identifies as a strength the author's refusal of both the absolute neutrality of science and the claim that all knowledge is arbitrary. However, the distinction between the two degrees of ideology depends on a judgment whose criteria could be further developed. The discussion is related to the film Oppenheimer, in which the technical ability to build the atomic bomb does not solve the moral problem of its use.}

\vspace{6pt}
\noindent\textit{\textbf{Keywords:} Science. Ideology. Ethics. Scientific discourse. Knowledge.}
\setlength{\parindent}{1.25cm}

\vspace{18pt}

% ======================================================
% RESENHA
% ======================================================
\begin{center}
\textbf{\large Resenha}
\end{center}
\vspace{6pt}

Esta é uma resenha do capítulo ``Ciência e ideologia'', escrito por Gérard Fourez e publicado no livro \textit{A construção das ciências: introdução à filosofia e à ética das ciências}. A edição usada foi traduzida por Luiz Paulo Rouanet e publicada pela Editora da Universidade Estadual Paulista, em 1995. O capítulo ocupa as páginas 179 a 193 e está dividido em sete partes, nas quais o autor apresenta os discursos ideológicos, discute sua crítica pela ciência e depois aplica essa discussão à própria atividade científica.

Gérard Fourez foi um intelectual belga, licenciado em Filosofia e Matemática e doutor em Física Teórica. Foi professor da Universidade de Namur, onde fundou um departamento dedicado ao estudo das relações entre a filosofia e o trabalho científico. Sua formação ajuda a situar o capítulo no campo da filosofia e da ética das ciências. Fourez não analisa a ciência apenas pelo funcionamento de seus métodos. Ele também considera os interesses, os valores e os contextos sociais presentes na produção do conhecimento (EDITORA UNESP, [s.d.]).

O assunto principal do capítulo é a relação entre ciência e ideologia. O problema pode ser formulado da seguinte maneira: se a ciência também é histórica e depende de escolhas, até que ponto ela pode criticar as ideologias sem agir como uma delas? Fourez responde que a ciência possui força crítica, mas não ocupa uma posição totalmente neutra. Ela pode testar afirmações específicas, comparar dados e mostrar incoerências. Mesmo assim, seus conceitos são construídos a partir de pontos de vista determinados. A ciência se torna mais problemática quando esconde essas condições e apresenta uma explicação parcial como verdade universal (FOUREZ, 1995, p. 179--189).

Na primeira parte, ``Discursos ideológicos e eficácia crítica da ciência'', o autor define como ideológico o discurso que parece descrever o mundo, mas serve principalmente para legitimar práticas, motivar pessoas e fortalecer a união de um grupo. Esse discurso seleciona o que será mostrado e oculta parte de seus critérios. Fourez cita afirmações sobre a suposta fragilidade das mulheres, a inteligência humana e a exploração dos países pobres. Em todos esses casos, palavras amplas são usadas como se tivessem um sentido evidente. O efeito ideológico não depende de uma mentira completa. Ele pode aparecer no destaque dado a algumas diferenças e no apagamento de semelhanças importantes (FOUREZ, 1995, p. 179--180).

O exemplo mais provocador dessa parte aproxima o contrato de trabalho da prostituição. Nos dois casos, uma pessoa permite que outra disponha de uma parte de sua intimidade ou criatividade em troca de dinheiro. Fourez não afirma que as duas práticas sejam iguais em todos os aspectos. Seu objetivo é mostrar que a maneira habitual de classificá-las faz uma semelhança desaparecer. O exemplo é eficiente porque obriga o leitor a perceber que até uma descrição aparentemente simples depende dos aspectos escolhidos por quem descreve.

Na segunda parte, ``Crítica da ideologia pela ciência'', Fourez explica como a ciência pode revelar os critérios escondidos em uma afirmação. Uma ideia geral, como dizer que o ser humano é mais inteligente do que o macaco, pode ser transformada em uma pergunta limitada e testável. Para isso, é necessário definir o que será considerado inteligência e escolher uma forma de medi-la. O resultado pode responder à questão específica, mas não prova a afirmação geral. O teste mede apenas a definição adotada. Assim, a tradução científica esclarece uma parte do problema, sem representar sua totalidade (FOUREZ, 1995, p. 180--184).

Essa tradução também apresenta um risco. Quando uma reclamação ampla é convertida em critérios técnicos, o sentimento que a originou pode ser eliminado. Fourez usa o exemplo de estudantes que dizem que os horários dos exames são ruins. Uma análise pode definir intervalos, quantidade de provas e tempo de estudo, mas ainda assim não expressar todo o descontentamento vivido por eles. A ciência continua sendo uma forma importante de crítica porque permite testes precisos. Porém, acontecimentos cotidianos e valores filosóficos, éticos ou religiosos também podem colocar uma ideologia em dúvida (FOUREZ, 1995, p. 182--184).

Na terceira parte, ``Incapacidade da ciência em esclarecer inteiramente as questões éticas'', esse limite fica mais claro. Conhecimentos científicos ajudam a descrever um embrião, uma doença ou o funcionamento do corpo, mas não decidem sozinhos quais direitos devem ser reconhecidos ou qual escolha deve ser feita. Essas decisões exigem valores. Fourez chama de redução o erro de tratar uma tradução científica parcial como explicação completa da experiência. A frase ``o amor é uma questão de hormônios'' é um exemplo: os hormônios participam dessa experiência, mas não esgotam o que ela significa para as pessoas (FOUREZ, 1995, p. 184--186).

As duas partes seguintes apresentam a diferença entre ideologias de primeiro e segundo graus. Uma ideologia de primeiro grau ainda permite reconhecer os vestígios de sua construção histórica. Nesse caso, o discurso admite seus limites e as decisões usadas para formar seus conceitos. Na ideologia de segundo grau, esses vestígios são apagados e uma escolha particular aparece como se fosse natural, objetiva e eterna. Fourez aplica a distinção à própria ciência. Se ela se apresenta como uma tecnologia intelectual histórica e limitada, permanece no primeiro grau. Quando promete respostas neutras e definitivas para qualquer problema, passa ao segundo (FOUREZ, 1995, p. 186--189).

O exemplo do transporte resume bem essa diferença. Dizer que o carro \textit{pode ser} uma resposta reconhece a existência de outras soluções e de um contexto. Dizer que ele \textit{é} a resposta oculta as escolhas históricas e econômicas que fizeram esse meio de transporte parecer necessário. O problema, então, não está em usar conceitos científicos, mas em esquecer que foram construídos para responder a perguntas delimitadas.

Na sexta parte, Fourez afirma que a transmissão de ideologias ocorre muitas vezes sem intenção. Uma pessoa pode acreditar que não possui uma visão preconceituosa e, mesmo assim, reproduzi-la em sua fala. A propaganda se diferencia porque utiliza conscientemente a ocultação em favor de um projeto político ou econômico. Como toda visão de mundo depende de critérios e de um meio social, não é possível eliminar qualquer ideologia. A responsabilidade ética está em buscar clareza sobre o que se transmite e decidir quais posições podem ser assumidas. O autor dirige essa preocupação especialmente a cientistas e professores, pois ambos ajudam a formar representações do mundo (FOUREZ, 1995, p. 189--191).

Na última parte, ``A ciência varia de acordo com o grupo social?'', o autor analisa a influência social sobre as prioridades científicas. A ciência moderna se desenvolveu junto à visão burguesa de domínio e exploração da natureza. A medicina, por exemplo, concentrou-se mais no tratamento individual do que na prevenção, em parte por ter sido organizada segundo as demandas das classes privilegiadas. Outros grupos podem formular perguntas e interesses diferentes. Isso não significa, contudo, que cada grupo possa inventar qualquer resposta. Fourez recorda o fracasso da biologia de Lyssenko na União Soviética para mostrar que a realidade impõe resistências. O conhecimento é uma construção social, mas não é uma criação arbitrária (FOUREZ, 1995, p. 191--193).

Esse ponto é um dos maiores méritos do capítulo. Fourez evita dois extremos. De um lado, rejeita a ideia de uma ciência completamente neutra, separada da sociedade. De outro, não conclui que todas as explicações tenham o mesmo valor. Projetos sociais influenciam as perguntas e os usos do conhecimento, mas os testes e a experiência limitam aquilo que pode ser sustentado. Essa posição é importante porque permite criticar a ciência sem negar sua capacidade de produzir conhecimentos confiáveis.

A discussão me fez pensar no filme \textit{Oppenheimer}, de Christopher Nolan. Os cientistas do Projeto Manhattan dominavam os conhecimentos necessários para construir a bomba atômica, mas esse domínio não respondia se ela deveria ser usada. A Física podia explicar o funcionamento da arma e estimar seus efeitos. A decisão sobre seu emprego envolvia interesses militares, objetivos políticos e julgamentos morais. O caso ilustra a tese de Fourez sem exigir que se considere a ciência inútil ou culpada por si mesma: o conhecimento técnico resolve uma parte do problema, enquanto a decisão prática ultrapassa essa parte (OPPENHEIMER, 2023).

Apesar de convincente, a distinção entre ideologia de primeiro e segundo graus deixa uma dificuldade. O próprio Fourez reconhece que ela não é objetiva e depende da decisão de quem avalia se um discurso oculta demais os seus critérios (FOUREZ, 1995, p. 188). Essa sinceridade fortalece o texto, mas também mostra uma limitação. O capítulo poderia oferecer procedimentos mais concretos para comparar essas avaliações. Sem critérios compartilhados, duas pessoas podem reconhecer a origem histórica de uma ideia e ainda discordar sobre o grau de ocultação presente nela. A diferença continua útil, porém exige discussão e justificativa, não uma classificação automática.

A principal contribuição do capítulo está em ensinar o leitor a perguntar como uma afirmação foi construída, quais interesses orientaram suas categorias e até onde seus resultados podem ser aplicados. A linguagem é acessível, e os exemplos tornam conceitos abstratos mais fáceis de acompanhar. Em alguns momentos, a quantidade de exemplos alonga a explicação, mas isso não prejudica a sequência do argumento: Fourez parte da crítica científica das ideologias, mostra os limites dessa crítica e termina examinando as condições sociais da própria ciência.

A leitura é recomendada principalmente a estudantes, professores e profissionais que trabalham com conhecimento científico e técnico. Também pode interessar a leitores que desejam compreender melhor o uso de expressões como neutralidade, objetividade e ideologia. Para estudantes de Ciência da Computação, o capítulo é pertinente porque lembra que uma solução tecnicamente possível não se torna, apenas por isso, socialmente desejável. O valor do texto está em defender uma ciência rigorosa, mas consciente de seus limites e de sua responsabilidade.

% ======================================================
% REFERÊNCIAS
% ======================================================
\begin{center}
\textbf{\large Referências}
\end{center}
\vspace{6pt}

\begin{singlespace}
\sloppy
\noindent\setlength{\parindent}{0pt}%
EDITORA UNESP. \textit{A construção das ciências: introdução à filosofia e à ética das ciências}. São Paulo: Fundação Editora da Unesp, [s.d.]. Disponível em: \url{https://editoraunesp.com.br/catalogo/9788571390836,a-construcao-das-ciencias}. Acesso em: 2 out. 2026.

\par\vspace{6pt}\noindent
FOUREZ, Gérard. Ciência e ideologia. In: FOUREZ, Gérard. \textit{A construção das ciências: introdução à filosofia e à ética das ciências}. Tradução de Luiz Paulo Rouanet. São Paulo: Editora da Universidade Estadual Paulista, 1995. cap. 7, p. 179--193.

\par\vspace{6pt}\noindent
OPPENHEIMER. Direção e roteiro: Christopher Nolan. Produção: Emma Thomas, Charles Roven e Christopher Nolan. Universal Pictures, 2023. 1 filme (181 min), son., color.
\end{singlespace}

\end{document}

